\documentclass{article}

% INFR11269 Machine Learning Systems — Coursework Report (ICLR 2027 style; do not change the format).
% Main body: six pages. References and Appendix do not count.

\usepackage{iclr2027_conference,times}

\usepackage{hyperref}
\usepackage{url}
\usepackage{booktabs}
\usepackage{graphicx}
\usepackage{amsmath}
\usepackage{amssymb}

%%%%% NEW MATH DEFINITIONS %%%%%

\usepackage{amsmath,amsfonts,bm}

% Mark sections of captions for referring to divisions of figures
\newcommand{\figleft}{{\em (Left)}}
\newcommand{\figcenter}{{\em (Center)}}
\newcommand{\figright}{{\em (Right)}}
\newcommand{\figtop}{{\em (Top)}}
\newcommand{\figbottom}{{\em (Bottom)}}
\newcommand{\captiona}{{\em (a)}}
\newcommand{\captionb}{{\em (b)}}
\newcommand{\captionc}{{\em (c)}}
\newcommand{\captiond}{{\em (d)}}

% Highlight a newly defined term
\newcommand{\newterm}[1]{{\bf #1}}


% Figure reference, lower-case.
\def\figref#1{figure~\ref{#1}}
% Figure reference, capital. For start of sentence
\def\Figref#1{Figure~\ref{#1}}
\def\twofigref#1#2{figures \ref{#1} and \ref{#2}}
\def\quadfigref#1#2#3#4{figures \ref{#1}, \ref{#2}, \ref{#3} and \ref{#4}}
% Section reference, lower-case.
\def\secref#1{section~\ref{#1}}
% Section reference, capital.
\def\Secref#1{Section~\ref{#1}}
% Reference to two sections.
\def\twosecrefs#1#2{sections \ref{#1} and \ref{#2}}
% Reference to three sections.
\def\secrefs#1#2#3{sections \ref{#1}, \ref{#2} and \ref{#3}}
% Reference to an equation, lower-case.
\def\eqref#1{equation~\ref{#1}}
% Reference to an equation, upper case
\def\Eqref#1{Equation~\ref{#1}}
% A raw reference to an equation---avoid using if possible
\def\plaineqref#1{\ref{#1}}
% Reference to a chapter, lower-case.
\def\chapref#1{chapter~\ref{#1}}
% Reference to an equation, upper case.
\def\Chapref#1{Chapter~\ref{#1}}
% Reference to a range of chapters
\def\rangechapref#1#2{chapters\ref{#1}--\ref{#2}}
% Reference to an algorithm, lower-case.
\def\algref#1{algorithm~\ref{#1}}
% Reference to an algorithm, upper case.
\def\Algref#1{Algorithm~\ref{#1}}
\def\twoalgref#1#2{algorithms \ref{#1} and \ref{#2}}
\def\Twoalgref#1#2{Algorithms \ref{#1} and \ref{#2}}
% Reference to a part, lower case
\def\partref#1{part~\ref{#1}}
% Reference to a part, upper case
\def\Partref#1{Part~\ref{#1}}
\def\twopartref#1#2{parts \ref{#1} and \ref{#2}}

\def\ceil#1{\lceil #1 \rceil}
\def\floor#1{\lfloor #1 \rfloor}
\def\1{\bm{1}}
\newcommand{\train}{\mathcal{D}}
\newcommand{\valid}{\mathcal{D_{\mathrm{valid}}}}
\newcommand{\test}{\mathcal{D_{\mathrm{test}}}}

\def\eps{{\epsilon}}


% Random variables
\def\reta{{\textnormal{$\eta$}}}
\def\ra{{\textnormal{a}}}
\def\rb{{\textnormal{b}}}
\def\rc{{\textnormal{c}}}
\def\rd{{\textnormal{d}}}
\def\re{{\textnormal{e}}}
\def\rf{{\textnormal{f}}}
\def\rg{{\textnormal{g}}}
\def\rh{{\textnormal{h}}}
\def\ri{{\textnormal{i}}}
\def\rj{{\textnormal{j}}}
\def\rk{{\textnormal{k}}}
\def\rl{{\textnormal{l}}}
% rm is already a command, just don't name any random variables m
\def\rn{{\textnormal{n}}}
\def\ro{{\textnormal{o}}}
\def\rp{{\textnormal{p}}}
\def\rq{{\textnormal{q}}}
\def\rr{{\textnormal{r}}}
\def\rs{{\textnormal{s}}}
\def\rt{{\textnormal{t}}}
\def\ru{{\textnormal{u}}}
\def\rv{{\textnormal{v}}}
\def\rw{{\textnormal{w}}}
\def\rx{{\textnormal{x}}}
\def\ry{{\textnormal{y}}}
\def\rz{{\textnormal{z}}}

% Random vectors
\def\rvepsilon{{\mathbf{\epsilon}}}
\def\rvtheta{{\mathbf{\theta}}}
\def\rva{{\mathbf{a}}}
\def\rvb{{\mathbf{b}}}
\def\rvc{{\mathbf{c}}}
\def\rvd{{\mathbf{d}}}
\def\rve{{\mathbf{e}}}
\def\rvf{{\mathbf{f}}}
\def\rvg{{\mathbf{g}}}
\def\rvh{{\mathbf{h}}}
\def\rvu{{\mathbf{i}}}
\def\rvj{{\mathbf{j}}}
\def\rvk{{\mathbf{k}}}
\def\rvl{{\mathbf{l}}}
\def\rvm{{\mathbf{m}}}
\def\rvn{{\mathbf{n}}}
\def\rvo{{\mathbf{o}}}
\def\rvp{{\mathbf{p}}}
\def\rvq{{\mathbf{q}}}
\def\rvr{{\mathbf{r}}}
\def\rvs{{\mathbf{s}}}
\def\rvt{{\mathbf{t}}}
\def\rvu{{\mathbf{u}}}
\def\rvv{{\mathbf{v}}}
\def\rvw{{\mathbf{w}}}
\def\rvx{{\mathbf{x}}}
\def\rvy{{\mathbf{y}}}
\def\rvz{{\mathbf{z}}}

% Elements of random vectors
\def\erva{{\textnormal{a}}}
\def\ervb{{\textnormal{b}}}
\def\ervc{{\textnormal{c}}}
\def\ervd{{\textnormal{d}}}
\def\erve{{\textnormal{e}}}
\def\ervf{{\textnormal{f}}}
\def\ervg{{\textnormal{g}}}
\def\ervh{{\textnormal{h}}}
\def\ervi{{\textnormal{i}}}
\def\ervj{{\textnormal{j}}}
\def\ervk{{\textnormal{k}}}
\def\ervl{{\textnormal{l}}}
\def\ervm{{\textnormal{m}}}
\def\ervn{{\textnormal{n}}}
\def\ervo{{\textnormal{o}}}
\def\ervp{{\textnormal{p}}}
\def\ervq{{\textnormal{q}}}
\def\ervr{{\textnormal{r}}}
\def\ervs{{\textnormal{s}}}
\def\ervt{{\textnormal{t}}}
\def\ervu{{\textnormal{u}}}
\def\ervv{{\textnormal{v}}}
\def\ervw{{\textnormal{w}}}
\def\ervx{{\textnormal{x}}}
\def\ervy{{\textnormal{y}}}
\def\ervz{{\textnormal{z}}}

% Random matrices
\def\rmA{{\mathbf{A}}}
\def\rmB{{\mathbf{B}}}
\def\rmC{{\mathbf{C}}}
\def\rmD{{\mathbf{D}}}
\def\rmE{{\mathbf{E}}}
\def\rmF{{\mathbf{F}}}
\def\rmG{{\mathbf{G}}}
\def\rmH{{\mathbf{H}}}
\def\rmI{{\mathbf{I}}}
\def\rmJ{{\mathbf{J}}}
\def\rmK{{\mathbf{K}}}
\def\rmL{{\mathbf{L}}}
\def\rmM{{\mathbf{M}}}
\def\rmN{{\mathbf{N}}}
\def\rmO{{\mathbf{O}}}
\def\rmP{{\mathbf{P}}}
\def\rmQ{{\mathbf{Q}}}
\def\rmR{{\mathbf{R}}}
\def\rmS{{\mathbf{S}}}
\def\rmT{{\mathbf{T}}}
\def\rmU{{\mathbf{U}}}
\def\rmV{{\mathbf{V}}}
\def\rmW{{\mathbf{W}}}
\def\rmX{{\mathbf{X}}}
\def\rmY{{\mathbf{Y}}}
\def\rmZ{{\mathbf{Z}}}

% Elements of random matrices
\def\ermA{{\textnormal{A}}}
\def\ermB{{\textnormal{B}}}
\def\ermC{{\textnormal{C}}}
\def\ermD{{\textnormal{D}}}
\def\ermE{{\textnormal{E}}}
\def\ermF{{\textnormal{F}}}
\def\ermG{{\textnormal{G}}}
\def\ermH{{\textnormal{H}}}
\def\ermI{{\textnormal{I}}}
\def\ermJ{{\textnormal{J}}}
\def\ermK{{\textnormal{K}}}
\def\ermL{{\textnormal{L}}}
\def\ermM{{\textnormal{M}}}
\def\ermN{{\textnormal{N}}}
\def\ermO{{\textnormal{O}}}
\def\ermP{{\textnormal{P}}}
\def\ermQ{{\textnormal{Q}}}
\def\ermR{{\textnormal{R}}}
\def\ermS{{\textnormal{S}}}
\def\ermT{{\textnormal{T}}}
\def\ermU{{\textnormal{U}}}
\def\ermV{{\textnormal{V}}}
\def\ermW{{\textnormal{W}}}
\def\ermX{{\textnormal{X}}}
\def\ermY{{\textnormal{Y}}}
\def\ermZ{{\textnormal{Z}}}

% Vectors
\def\vzero{{\bm{0}}}
\def\vone{{\bm{1}}}
\def\vmu{{\bm{\mu}}}
\def\vtheta{{\bm{\theta}}}
\def\va{{\bm{a}}}
\def\vb{{\bm{b}}}
\def\vc{{\bm{c}}}
\def\vd{{\bm{d}}}
\def\ve{{\bm{e}}}
\def\vf{{\bm{f}}}
\def\vg{{\bm{g}}}
\def\vh{{\bm{h}}}
\def\vi{{\bm{i}}}
\def\vj{{\bm{j}}}
\def\vk{{\bm{k}}}
\def\vl{{\bm{l}}}
\def\vm{{\bm{m}}}
\def\vn{{\bm{n}}}
\def\vo{{\bm{o}}}
\def\vp{{\bm{p}}}
\def\vq{{\bm{q}}}
\def\vr{{\bm{r}}}
\def\vs{{\bm{s}}}
\def\vt{{\bm{t}}}
\def\vu{{\bm{u}}}
\def\vv{{\bm{v}}}
\def\vw{{\bm{w}}}
\def\vx{{\bm{x}}}
\def\vy{{\bm{y}}}
\def\vz{{\bm{z}}}

% Elements of vectors
\def\evalpha{{\alpha}}
\def\evbeta{{\beta}}
\def\evepsilon{{\epsilon}}
\def\evlambda{{\lambda}}
\def\evomega{{\omega}}
\def\evmu{{\mu}}
\def\evpsi{{\psi}}
\def\evsigma{{\sigma}}
\def\evtheta{{\theta}}
\def\eva{{a}}
\def\evb{{b}}
\def\evc{{c}}
\def\evd{{d}}
\def\eve{{e}}
\def\evf{{f}}
\def\evg{{g}}
\def\evh{{h}}
\def\evi{{i}}
\def\evj{{j}}
\def\evk{{k}}
\def\evl{{l}}
\def\evm{{m}}
\def\evn{{n}}
\def\evo{{o}}
\def\evp{{p}}
\def\evq{{q}}
\def\evr{{r}}
\def\evs{{s}}
\def\evt{{t}}
\def\evu{{u}}
\def\evv{{v}}
\def\evw{{w}}
\def\evx{{x}}
\def\evy{{y}}
\def\evz{{z}}

% Matrix
\def\mA{{\bm{A}}}
\def\mB{{\bm{B}}}
\def\mC{{\bm{C}}}
\def\mD{{\bm{D}}}
\def\mE{{\bm{E}}}
\def\mF{{\bm{F}}}
\def\mG{{\bm{G}}}
\def\mH{{\bm{H}}}
\def\mI{{\bm{I}}}
\def\mJ{{\bm{J}}}
\def\mK{{\bm{K}}}
\def\mL{{\bm{L}}}
\def\mM{{\bm{M}}}
\def\mN{{\bm{N}}}
\def\mO{{\bm{O}}}
\def\mP{{\bm{P}}}
\def\mQ{{\bm{Q}}}
\def\mR{{\bm{R}}}
\def\mS{{\bm{S}}}
\def\mT{{\bm{T}}}
\def\mU{{\bm{U}}}
\def\mV{{\bm{V}}}
\def\mW{{\bm{W}}}
\def\mX{{\bm{X}}}
\def\mY{{\bm{Y}}}
\def\mZ{{\bm{Z}}}
\def\mBeta{{\bm{\beta}}}
\def\mPhi{{\bm{\Phi}}}
\def\mLambda{{\bm{\Lambda}}}
\def\mSigma{{\bm{\Sigma}}}

% Tensor
\DeclareMathAlphabet{\mathsfit}{\encodingdefault}{\sfdefault}{m}{sl}
\SetMathAlphabet{\mathsfit}{bold}{\encodingdefault}{\sfdefault}{bx}{n}
\newcommand{\tens}[1]{\bm{\mathsfit{#1}}}
\def\tA{{\tens{A}}}
\def\tB{{\tens{B}}}
\def\tC{{\tens{C}}}
\def\tD{{\tens{D}}}
\def\tE{{\tens{E}}}
\def\tF{{\tens{F}}}
\def\tG{{\tens{G}}}
\def\tH{{\tens{H}}}
\def\tI{{\tens{I}}}
\def\tJ{{\tens{J}}}
\def\tK{{\tens{K}}}
\def\tL{{\tens{L}}}
\def\tM{{\tens{M}}}
\def\tN{{\tens{N}}}
\def\tO{{\tens{O}}}
\def\tP{{\tens{P}}}
\def\tQ{{\tens{Q}}}
\def\tR{{\tens{R}}}
\def\tS{{\tens{S}}}
\def\tT{{\tens{T}}}
\def\tU{{\tens{U}}}
\def\tV{{\tens{V}}}
\def\tW{{\tens{W}}}
\def\tX{{\tens{X}}}
\def\tY{{\tens{Y}}}
\def\tZ{{\tens{Z}}}


% Graph
\def\gA{{\mathcal{A}}}
\def\gB{{\mathcal{B}}}
\def\gC{{\mathcal{C}}}
\def\gD{{\mathcal{D}}}
\def\gE{{\mathcal{E}}}
\def\gF{{\mathcal{F}}}
\def\gG{{\mathcal{G}}}
\def\gH{{\mathcal{H}}}
\def\gI{{\mathcal{I}}}
\def\gJ{{\mathcal{J}}}
\def\gK{{\mathcal{K}}}
\def\gL{{\mathcal{L}}}
\def\gM{{\mathcal{M}}}
\def\gN{{\mathcal{N}}}
\def\gO{{\mathcal{O}}}
\def\gP{{\mathcal{P}}}
\def\gQ{{\mathcal{Q}}}
\def\gR{{\mathcal{R}}}
\def\gS{{\mathcal{S}}}
\def\gT{{\mathcal{T}}}
\def\gU{{\mathcal{U}}}
\def\gV{{\mathcal{V}}}
\def\gW{{\mathcal{W}}}
\def\gX{{\mathcal{X}}}
\def\gY{{\mathcal{Y}}}
\def\gZ{{\mathcal{Z}}}

% Sets
\def\sA{{\mathbb{A}}}
\def\sB{{\mathbb{B}}}
\def\sC{{\mathbb{C}}}
\def\sD{{\mathbb{D}}}
% Don't use a set called E, because this would be the same as our symbol
% for expectation.
\def\sF{{\mathbb{F}}}
\def\sG{{\mathbb{G}}}
\def\sH{{\mathbb{H}}}
\def\sI{{\mathbb{I}}}
\def\sJ{{\mathbb{J}}}
\def\sK{{\mathbb{K}}}
\def\sL{{\mathbb{L}}}
\def\sM{{\mathbb{M}}}
\def\sN{{\mathbb{N}}}
\def\sO{{\mathbb{O}}}
\def\sP{{\mathbb{P}}}
\def\sQ{{\mathbb{Q}}}
\def\sR{{\mathbb{R}}}
\def\sS{{\mathbb{S}}}
\def\sT{{\mathbb{T}}}
\def\sU{{\mathbb{U}}}
\def\sV{{\mathbb{V}}}
\def\sW{{\mathbb{W}}}
\def\sX{{\mathbb{X}}}
\def\sY{{\mathbb{Y}}}
\def\sZ{{\mathbb{Z}}}

% Entries of a matrix
\def\emLambda{{\Lambda}}
\def\emA{{A}}
\def\emB{{B}}
\def\emC{{C}}
\def\emD{{D}}
\def\emE{{E}}
\def\emF{{F}}
\def\emG{{G}}
\def\emH{{H}}
\def\emI{{I}}
\def\emJ{{J}}
\def\emK{{K}}
\def\emL{{L}}
\def\emM{{M}}
\def\emN{{N}}
\def\emO{{O}}
\def\emP{{P}}
\def\emQ{{Q}}
\def\emR{{R}}
\def\emS{{S}}
\def\emT{{T}}
\def\emU{{U}}
\def\emV{{V}}
\def\emW{{W}}
\def\emX{{X}}
\def\emY{{Y}}
\def\emZ{{Z}}
\def\emSigma{{\Sigma}}

% entries of a tensor
% Same font as tensor, without \bm wrapper
\newcommand{\etens}[1]{\mathsfit{#1}}
\def\etLambda{{\etens{\Lambda}}}
\def\etA{{\etens{A}}}
\def\etB{{\etens{B}}}
\def\etC{{\etens{C}}}
\def\etD{{\etens{D}}}
\def\etE{{\etens{E}}}
\def\etF{{\etens{F}}}
\def\etG{{\etens{G}}}
\def\etH{{\etens{H}}}
\def\etI{{\etens{I}}}
\def\etJ{{\etens{J}}}
\def\etK{{\etens{K}}}
\def\etL{{\etens{L}}}
\def\etM{{\etens{M}}}
\def\etN{{\etens{N}}}
\def\etO{{\etens{O}}}
\def\etP{{\etens{P}}}
\def\etQ{{\etens{Q}}}
\def\etR{{\etens{R}}}
\def\etS{{\etens{S}}}
\def\etT{{\etens{T}}}
\def\etU{{\etens{U}}}
\def\etV{{\etens{V}}}
\def\etW{{\etens{W}}}
\def\etX{{\etens{X}}}
\def\etY{{\etens{Y}}}
\def\etZ{{\etens{Z}}}

% The true underlying data generating distribution
\newcommand{\pdata}{p_{\rm{data}}}
% The empirical distribution defined by the training set
\newcommand{\ptrain}{\hat{p}_{\rm{data}}}
\newcommand{\Ptrain}{\hat{P}_{\rm{data}}}
% The model distribution
\newcommand{\pmodel}{p_{\rm{model}}}
\newcommand{\Pmodel}{P_{\rm{model}}}
\newcommand{\ptildemodel}{\tilde{p}_{\rm{model}}}
% Stochastic autoencoder distributions
\newcommand{\pencode}{p_{\rm{encoder}}}
\newcommand{\pdecode}{p_{\rm{decoder}}}
\newcommand{\precons}{p_{\rm{reconstruct}}}

\newcommand{\laplace}{\mathrm{Laplace}} % Laplace distribution

\newcommand{\E}{\mathbb{E}}
\newcommand{\Ls}{\mathcal{L}}
\newcommand{\R}{\mathbb{R}}
\newcommand{\emp}{\tilde{p}}
\newcommand{\lr}{\alpha}
\newcommand{\reg}{\lambda}
\newcommand{\rect}{\mathrm{rectifier}}
\newcommand{\softmax}{\mathrm{softmax}}
\newcommand{\sigmoid}{\sigma}
\newcommand{\softplus}{\zeta}
\newcommand{\KL}{D_{\mathrm{KL}}}
\newcommand{\Var}{\mathrm{Var}}
\newcommand{\standarderror}{\mathrm{SE}}
\newcommand{\Cov}{\mathrm{Cov}}
% Wolfram Mathworld says $L^2$ is for function spaces and $\ell^2$ is for vectors
% But then they seem to use $L^2$ for vectors throughout the site, and so does
% wikipedia.
\newcommand{\normlzero}{L^0}
\newcommand{\normlone}{L^1}
\newcommand{\normltwo}{L^2}
\newcommand{\normlp}{L^p}
\newcommand{\normmax}{L^\infty}

\newcommand{\parents}{Pa} % See usage in notation.tex. Chosen to match Daphne's book.

\DeclareMathOperator*{\argmax}{arg\,max}
\DeclareMathOperator*{\argmin}{arg\,min}

\DeclareMathOperator{\sign}{sign}
\DeclareMathOperator{\Tr}{Tr}
\let\ab\allowbreak


\title{Where the Time and the Evidence Go:\\CUDA Kernels and Retrieval Depth in a Spoken-Question RAG System}

% Keep this line unchanged.
\author{Anonymous coursework submission}

% Keep this COMMENTED OUT.
% \iclrfinalcopy

\begin{document}

\maketitle

% Course-specific template fingerprint. Keep this line unchanged.
\lhead{Under review as a report for UoE 26/27 Sem1 INFR11269: Machine Learning Systems}

% ============================================================
\section{Introduction}
\label{sec:introduction}

A spoken question travels through four stages before it is answered: speech recognition (ASR), a sentence encoder, dense retrieval over a 24{,}000-record finance corpus, and a large language model (LLM) that writes the answer from the retrieved records \citep{lewis2020rag}. We worked on both ends of this chain. For ASR we implemented the four GPU kernels the course model needs (RoPE, block reductions, tiled matrix multiply, tiled attention) and measured where the time goes. For retrieval-augmented generation (RAG) we asked one question: \emph{how much retrieved evidence should the generator receive when the question comes from ASR?} Retrieval depth $k$ is the main lever on generation cost, because every extra record is about 180 words of prompt, and it is also the obvious lever on answer quality, because more records raise the chance that the right one is present.

\textbf{Baselines.} ASR: the course model with PyTorch operators, on the same GPU node and in the same job as ours. RAG: the supplied pipeline with $k=5$.
\textbf{Our changes.} The four kernels with a profile-guided choice of tile size and queries per block; a depth study at $k\in\{1,5,20\}$ with paired, judged answers; and an adaptive depth policy that we designed from the depth study's mechanism.

\textbf{Main findings.} (1) Our kernels are correct on every course test and transcribe the test recording exactly, but the full model is $2.45\times$ slower than PyTorch; 86\% of kernel time is the matrix multiply, which reaches 16--31\% of cuBLAS. Choosing tile 32 and 16 attention warps from ablations made it 17\% faster than the handout settings. (2) Recall rises steeply with $k$, but judged answer quality shows no measurable change: over 55 judged queries from two splits, $k=1$, 5 and 20 score 0.267, 0.294 and 0.273 while the prompt grows $15\times$. Grouping queries by the rank of their first relevant record suggests why: extra records significantly help queries whose evidence is ranked low ($k=1\to5$), while queries whose top record is already right lose slightly with depth on both splits (not significant), so in the average the two largely net out. (3) The adaptive policy does not beat a fixed depth; even a perfect gate would gain little here, because many answers score 0 at every depth. The answers to the four required task questions are collected in Section~\ref{sec:questions}.

% ============================================================
\section{System and Technical Approach}
\label{sec:approach}

\textbf{Supplied.} The ASR model (GLM-ASR, an audio encoder of 32 layers over 1500 frames and a text decoder with grouped-query attention, GQA), its tests and benchmark; the RAG framework (data, a 384-dimensional encoder, exact cosine retrieval, a context builder that concatenates records, an LLM generator and an LLM judge). Generator and judge are the same model (EuroLLM-22B-Instruct) through the university's ELM service, at temperature 0.1 and 0.0. \textbf{Ours.} The five task files, an ablation harness that rebuilds each kernel with one or two compile-time flags, and the RAG experiment scripts (Appendix~\ref{app:code}). No supplied file was changed; a checksum guard re-checks all 55 supplied files each time our files are synchronised into the course checkout.

\subsection{ASR kernels}
\textbf{RoPE.} One launch rotates the query and key heads together: each thread owns one (position, head, pair) element, computes $\sin$ and $\cos$ of position $\times$ inverse frequency, and rotates the pair \citep{su2024roformer}. The kernel is memory-bound, so the design goal is one read and one write per element.

\textbf{Reductions.} Warp sums and maxima use register shuffles; a block reduction combines the warps through shared memory, with $-\infty$ (not 0) for missing warps in the maximum. RMSNorm and LayerNorm (the latter, like the block sum, supplied as a worked example) reduce one row per block. Softmax uses the online (maximum, sum) merge \citep{milakov2018online}: each thread keeps a running maximum and a sum rescaled to it, the block merges these pairs, and a second pass writes $\exp(x-m)/s$. This reads the row twice instead of three times.

\textbf{Tiled GEMM.} Each block computes a TILE$\times$TILE output tile, staging tiles of the input and of the weight matrix in shared memory so that each loaded value is reused TILE times. The weight tile is stored with one padding column so that the 32 lanes of a warp read 32 different shared-memory banks, and bias and GELU are applied in the epilogue, saving a separate kernel's read and write.

\textbf{Attention.} One warp per query, ATTN\_WARPS queries per block; the block loads keys and values into shared memory one tile at a time and every warp keeps a running maximum, normaliser and output accumulator (online softmax, as in FlashAttention \citep{dao2022flashattention}), so the $1500\times1500$ score matrix is never stored. GQA maps query head $h$ to key head $h/(H_q/H_k)$; causal masking stops each query at its own last key, while all warps still execute the same number of tile iterations so that block barriers stay valid.

\textbf{Design choice: tune only what the profile says matters.} The per-operation profile (Section~\ref{sec:results}, Appendix~\ref{app:ablations}) shows the matrix multiply and attention take over 99\% of kernel time, so we ran the guides' experiments for all four tasks but chose final settings only for those two: TILE~32 (handout 16) and ATTN\_WARPS~16 (handout 4).

\subsection{RAG investigation}
\textbf{Depth study.} We vary only $k$ and keep everything else fixed: ASR ran once on the GPU and its transcripts were cached, so every $k$ sees identical text; prompt, model, temperature and judge are the supplied ones. \textbf{Adaptive depth.} The depth study suggested that the best $k$ depends on whether the top record is already right (Section~\ref{sec:mechanism}). The policy uses the retriever's own scores, which are free: $k=1$ if the margin between the top two cosine scores $s_1-s_2$ is at least $\tau$, else $k=20$. To avoid tuning on the outcome, the protocol was fixed in advance: choose the feature ($s_1$ or the margin) by AUC for ``rank-1 record is relevant'' and $\tau$ by Youden's $J$, using only the practice split's relevance labels; then freeze both and test on the held-out dev split.

% ============================================================
\section{Experimental Methodology}
\label{sec:methodology}

\textbf{Hardware.} ASR on an NVIDIA RTX 2080 Ti (sm\_75, 68 SMs, 616\,GB/s peak), pinned by Slurm resource type on a teaching cluster, CUDA 12.8. Retrieval and encoding on a laptop CPU; generation and judging on the remote ELM service.

\textbf{ASR measurements.} Correctness: the course tests compare each kernel with a CPU or cuBLAS reference on edge cases (odd sizes, inputs of $\pm100$, all-negative rows, one new query) and real model shapes; compute-sanitizer memcheck and racecheck on the small test cases; the full model must reproduce the reference transcript of the test recording. Cost: milliseconds per transcription of the 3.5\,s test recording (13 tokens), the mean of 3 runs, model loading and feature extraction excluded; per-operation GPU time from the course profiler; GB/s and GFLOP/s from the tests. PyTorch's time varies by node (806 to 1215\,ms on three nodes of the same GPU type), which we suspect comes from host-side launch overhead of its many small operators but did not test; ours does not (repeats of one configuration on different nodes agree within 0.3--1.1\%). We therefore compare with PyTorch only within one job on one node, and treat differences of about 1\% or less as noise.

\textbf{RAG data and metrics.} 350 practice and 125 dev spoken queries, with relevance labels (qrels) but no reference transcripts. Retrieval: Recall@$k$ and hit@$k$ (at least one relevant record in the top $k$) over all queries. Answers: the supplied judge scores correctness, grounding and completeness from 0 to 4 against the gold records; we report the normalised sum (0--1). We judged the first 30 practice and the first 28 dev queries at every $k$, paired (the same queries at every depth); the judge's output could not be parsed for 3 of 58 queries, which are excluded at every $k$ (practice $n=28$, dev $n=27$). Cost: prompt tokens returned by the API, generation latency (including network and service queueing), context words. Intervals are 95\% bootstrap intervals over queries (2000 resamples per split, 5000 for the pooled set), and differences between depths are paired. The generator (temperature 0.1) and the judge (temperature 0) ran once per answer; a second judge pass over 56 stored answers measures the judge's agreement with itself (Section~\ref{sec:conclusion}). The transcripts were produced once with the PyTorch ASR path on the GPU (mean 1555\,ms per query, model loading excluded) and cached, so RAG timings exclude ASR. Generation latency depends on the shared service's load, and the three depths ran one after another, so latency ratios across $k$ are approximate.

% ============================================================
\section{Results and Analysis}
\label{sec:results}

\subsection{ASR: correct, 2.45$\times$ slower than PyTorch, limited by the matrix multiply}

All four kernels pass every course test, both sanitizers report zero errors, and the full model on our kernels transcribes the test recording exactly. Table~\ref{tab:asr} gives the cost. Switching on our kernels one task at a time on one node shows that RoPE and the reductions make the model \emph{faster} than PyTorch (806 to 705\,ms), because PyTorch builds RoPE and RMSNorm from chains of small operators, while our tiled GEMM makes it $4.3\times$ slower. In the profile of our submitted configuration, the matrix multiply is 86.1\% of kernel time and attention 13.3\%; RoPE, both norms and softmax together are 0.5\%. On the PyTorch path the profiler attributes 47.6\% to matrix multiplies, 22.5\% to attention and 25.0\% to RoPE (handout run, another node); per-operation timing overstates PyTorch's many small operators (the profile sums to 1193\,ms against 806\,ms wall clock), and replacing RoPE saved 79\,ms end to end. Per generated token ours takes 207.6\,ms against 84.3\,ms, and our peak GPU memory is lower (8.56 vs 8.88\,GiB), consistent with attention never storing the score matrix. The memory-bound kernels are close to the hardware limit on the large cases: RoPE reaches 80\% of the 616\,GB/s peak and LayerNorm 62\%; RMSNorm on the 64-token decoder case (27\%) and the single-row vocabulary softmax (1\%, one block on one of 68 SMs) are too small to fill the GPU.

\begin{table}[t]
\caption{ASR cost on an RTX 2080 Ti, ms per transcription of the test recording (mean of 3 runs; $\pm$ standard deviation on the left). Left: the first three rows are one job on node A (PyTorch, ours, PyTorch); rows marked $\dagger$ are one job on node C. Right: kernels switched on one task at a time with the handout settings (one job, node B). All runs transcribe the test recording exactly.}
\label{tab:asr}
\begin{center}
\footnotesize
\setlength{\tabcolsep}{4pt}
\resizebox{\textwidth}{!}{%
\begin{tabular}{lr@{\hspace{1.5em}}lr}
\toprule
\multicolumn{2}{c}{Node A, one job ($\dagger$: node C)} & \multicolumn{2}{c}{Incremental, node B (handout settings)} \\
\midrule
PyTorch only & 1096.2 $\pm$ 1.8 & PyTorch only & 806.3 \\
Ours, submitted (tile 32, 16 warps) & 2698.8 $\pm$ 7.9 & + RoPE (Task 1) & 727.1 \\
PyTorch only, repeated & 1103.8 $\pm$ 4.9 & + reductions (Task 2) & 705.2 \\
\cmidrule{1-2}
Ours, handout (tile 16, 4 warps)$^\dagger$ & 3294.0 $\pm$ 8.3 & + GEMM (Task 3) & 2998.2 \\
Ours, tile 32 only$^\dagger$ & 2876.6 $\pm$ 9.6 & + attention (Task 4) & 3288.3 \\
Ours, 16 warps only$^\dagger$ & 3148.9 $\pm$ 6.9 & & \\
\bottomrule
\end{tabular}}
\end{center}
\footnotesize{$^\dagger$ One job on node C. The submitted settings took 2698.8 and 2711.3\,ms on node A (two jobs) and 2728.1\,ms on node C, within 1.1\%.}
\end{table}

\textbf{Design choices that mattered.} Tile 32 alone saves 13\% end to end and 16 attention warps 4\%; together 17.2\% (3294 to 2728\,ms on one node), almost exactly the sum, because they act on different kernels. Two layout details matter more than any tuning constant: without the padding column of the weight tile the large multiplies are $1.7$--$2.3\times$ slower at tile 16 and $4.5\times$ slower at tile 32 (every lane hits the same bank), and without the $+1$ padding of the attention key tile the encoder case is $4.7\times$ slower. The textbook three-pass softmax is 27\% slower than our two-pass version on the 59{,}264-wide vocabulary row.

\textbf{The gap to PyTorch is the matrix multiply} (Table~\ref{tab:gemm}). Tiling gives about $9\times$ over the naive kernel on every multi-row shape, but cuBLAS is another $3$--$6\times$ faster. A rough estimate (not a measurement) shows why: with one output per thread, every multiply-add needs its own shared-memory load of $w$ (plus a broadcast of $x$), and an SM serves one warp-wide shared load per cycle while it can issue two warp-wide multiply-adds, so such a kernel stays well below half the FP32 peak; cuBLAS computes many outputs per thread from registers. The single-row \texttt{lm\_head} is the exception: a $32\times32$ tile keeps only one of its 32 rows busy, so tile 32 is slower than naive there, and tile 8 was fastest in our sweep. One constant cannot suit every shape.

\begin{table}[t]
\caption{Matrix multiply throughput (GFLOP/s) on model shapes, from the course test on the hand-in run; ours uses tile 32.}
\label{tab:gemm}
\begin{center}
\footnotesize
\setlength{\tabcolsep}{4pt}
\begin{tabular}{lrrrrr}
\toprule
Shape ($M\times N\times K$) & Naive & Ours & cuBLAS & Ours/naive & Ours/cuBLAS \\
\midrule
Dec.\ q\_proj ($64\times2048\times2048$) & 153 & 1462 & 4647 & 9.6$\times$ & 31\% \\
Dec.\ down\_proj ($64\times2048\times6144$) & 193 & 1825 & 8143 & 9.5$\times$ & 22\% \\
Enc.\ fc1 + GELU ($1500\times5120\times1280$) & 233 & 2107 & 13363 & 9.0$\times$ & 16\% \\
Enc.\ conv2, im2col ($1500\times1280\times3840$) & 232 & 2130 & 13340 & 9.2$\times$ & 16\% \\
lm\_head, one token ($1\times59264\times2048$) & 111 & 63 & 291 & 0.57$\times$ & 22\% \\
\bottomrule
\end{tabular}
\end{center}
\end{table}

\subsection{Answers to the task questions}
\label{sec:questions}

\textbf{Task 1, Q1 (bytes and bandwidth, audio-encoder case).} The case rotates 1500 positions $\times$ 20 heads $\times$ 16 pairs (32 of 64 dimensions) for both $q$ and $k$: 960{,}000 pairs. Each pair reads two floats and writes two, 16 bytes, so the kernel moves 15.36\,MB. It takes 0.031\,ms, i.e.\ 491.8\,GB/s, 80\% of the 616\,GB/s peak. This is why the block size barely matters for RoPE between 64 and 512 threads: the kernel is already close to the memory bandwidth limit.

\textbf{Task 2, Q3 (reads of $x$ in softmax).} Ours reads each element twice: the first pass finds the maximum and the rescaled sum together with the online merge, and the second reads $x$ again to write the output. The textbook version reads it three times (maximum, sum, write). On the vocabulary row the extra pass costs 27\% (0.150 vs 0.118\,ms); on short rows the two are equal at the timer's resolution, because those calls (0.004--0.009\,ms) sit near the fixed cost of a kernel launch, which hides the extra pass; the 237\,KB vocabulary row also fits in the 5.5\,MB L2 cache, so cache residency alone does not separate the cases.

\textbf{Task 3, Q1 (global loads, $M=N=K=1024$).} The naive kernel loads one row of $x$ and one column of $w$ per output, $2K$ floats, so $2\times1024^3 = 2{,}147{,}483{,}648$ in total. With tiles of width 32 each staged value serves 32 outputs, so $2\times1024^3/32 = 67{,}108{,}864$, $32\times$ fewer. The measured gain is about $9\times$ (Table~\ref{tab:gemm}), because caches already absorb part of the naive kernel's repeated loads and ours is then limited by shared-memory reads and one output per thread; computing several outputs per thread from registers, as cuBLAS does \citep{volkov2008benchmarking}, is the next step.

\textbf{Task 4, Q4 (key loads and the limit of ATTN\_WARPS).} A block loads every key of its head once, and a head of 1500 queries has $1500/4=375$ blocks with 4 warps, so each key is loaded 375 times per head per layer (94 times with our 16 warps). Measured on the encoder case, 1, 2, 4, 8 and 16 warps take 103.6, 30.1, 11.5, 8.3 and 7.6\,ms: $9\times$ from 1 to 4 warps for $4\times$ fewer key loads. A likely reason, computed from the kernel's shared memory but not measured, is occupancy: a block needs about 16.5\,KB plus 256\,B per warp, so about three blocks fit on an SM, giving about 3, 6, 12, 24 and 32 resident warps for 1 to 16 warps per block; small blocks lose both key reuse and latency hiding. The gain shrinks (27\% from 4 to 8, 9\% from 8 to 16) once occupancy saturates and key loading stops being the bottleneck. The limits are 32 warps per block (1024 threads), shared memory and registers per block, and that fewer, larger blocks leave SMs idle at the end of the grid.

\subsection{RAG: more evidence, no measurable gain in answers}

Retrieval itself is cheap and does not depend on $k$ (under 1\,ms for an exact search over 24{,}000 records; encoding the query takes 21--54\,ms in 5 timed queries), but coverage rises steeply with depth (Table~\ref{tab:rag}): on practice, Recall@$k$ grows from 0.229 at $k=1$ to 0.579 at $k=20$, and the dev split agrees. Even at $k=20$, a quarter of queries have no relevant record at all (hit@20 = 0.737 practice, 0.768 dev).

\begin{table}[t]
\caption{Retrieval depth: quality and cost. Recall and hit@$k$ over all queries; judge score (0--1, 95\% bootstrap CI) on the judged subset, the same queries at every $k$; mean prompt tokens and median generation latency on the judged subset.}
\label{tab:rag}
\begin{center}
\footnotesize
\setlength{\tabcolsep}{4pt}
\begin{tabular}{lccc}
\toprule
 & $k=1$ & $k=5$ (baseline) & $k=20$ \\
\midrule
Recall@$k$, practice / dev ($n$=350 / 125) & 0.229 / 0.209 & 0.421 / 0.379 & 0.579 / 0.580 \\
hit@$k$, practice / dev & 0.411 / 0.376 & 0.591 / 0.584 & 0.737 / 0.768 \\
\midrule
Judge score, practice ($n$=28) & 0.250 (0.13--0.38) & 0.268 (0.16--0.39) & 0.292 (0.18--0.43) \\
Judge score, dev ($n$=27, held out) & 0.284 (0.16--0.41) & 0.321 (0.20--0.44) & 0.253 (0.12--0.39) \\
Judge score, pooled ($n$=55) & 0.267 & 0.294 & 0.273 \\
\midrule
Prompt tokens, practice / dev & 305 / 315 & 1249 / 1022 & 4730 / 4351 \\
Generation, median ms, practice / dev & 1803 / 2222 & 3782 / 5271 & 4303 / 7334 \\
\bottomrule
\end{tabular}
\end{center}
\end{table}

Judged answer quality does not follow coverage. Pooled over both splits, the paired differences are $k5-k1=+0.027$ (95\% CI $-0.055$ to $+0.108$), $k20-k1=+0.006$ ($-0.083$ to $+0.097$) and $k20-k5=-0.021$ ($-0.092$ to $+0.052$): all within $\pm0.03$, none significant, while the prompt grows $15\times$ from $k=1$ to $k=20$ and median generation time $2.4$--$3.3\times$. An earlier batch of 15 queries had suggested a clear gain for $k=20$ (0.361 vs 0.178); with the next 13 queries under identical settings the gain vanished, so we withdrew that claim. The three judge dimensions move with the total (dev correctness 1.15, 1.37 and 1.04 of 4 at $k=1,5,20$; Appendix~\ref{app:rag}). The study can rule out large effects, not small ones: the 95\% intervals of the differences are about $\pm0.08$ wide, so it excludes gains above about 0.11 but would likely miss effects below about 0.10--0.13 (80\% power).

\subsection{Why the averages are flat: coverage against dilution}
\label{sec:mechanism}

Grouping the 28 judged practice queries by the rank of their first relevant record suggests why the average is flat (a post hoc analysis). Queries whose evidence is at rank 1 ($n=11$) score best at $k=1$ and lose as records are added (0.447, 0.394, 0.386 at $k=1,5,20$), consistent with extra records acting as distractors \citep{liu2024lost}. Queries whose evidence is at ranks 2--5 ($n=8$) gain (0.208, 0.250, 0.344), because depth brings the evidence in. Queries with no relevant record in the top 20 ($n=5$) score 0 at every depth. In the average these effects largely net out. These groups were formed after seeing the practice scores and are small. On the 27 dev queries, which played no part in forming them, rank-1 queries again lose with depth (0.528, 0.458, 0.410; $n=12$) and queries with evidence at ranks 2--20 gain from $k=1$ to $k=5$ (0.131 to 0.381; $n=7$), though not further at $k=20$ (0.202). Pooled over both splits (paired bootstrap), the coverage effect is significant (ranks 2--20, $n=19$: $k5-k1=+0.149$, CI $+0.035$ to $+0.263$) and larger than in the rank-1 group (difference $+0.211$, CI $+0.024$ to $+0.399$), while the loss at rank 1 is not ($n=23$: $k20-k1=-0.091$, CI $-0.239$ to $+0.069$). The next experiment tests whether this can be exploited.

\begin{table}[t]
\caption{Left: judge score by the rank of the first relevant record (28 judged practice queries). Right: adaptive depth on the 27 held-out dev queries, paired with fixed $k=5$ (95\% CI); the oracle gives $k=1$ exactly when the rank-1 record is relevant.}
\label{tab:mech}
\begin{center}
\footnotesize
\setlength{\tabcolsep}{4pt}
\begin{tabular}{lcc@{\hspace{1.5em}}lccc}
\toprule
First relevant at & $n$ & $k$ = 1 / 5 / 20 & Policy (dev) & Score & vs $k=5$ & Tokens \\
\midrule
rank 1 & 11 & 0.447 / 0.394 / 0.386 & fixed $k=5$ & 0.321 & -- & 1022 \\
ranks 2--5 & 8 & 0.208 / 0.250 / 0.344 & fixed $k=20$ & 0.253 & $-0.068$ & 4351 \\
ranks 6--20 & 4 & 0.104 / 0.292 / 0.292 & adaptive & 0.272 & $-0.049$ ($-0.12$, $+0.02$) & 3367 \\
not in top 20 & 5 & 0 / 0 / 0 & oracle & 0.306 & $-0.015$ ($-0.09$, $+0.06$) & 2419 \\
\bottomrule
\end{tabular}
\end{center}
\end{table}

\textbf{Adaptive depth.} If the two effects can be separated, a policy that gives $k=1$ to queries whose top record is right and $k=20$ to the rest should beat any fixed $k$. On practice labels the margin feature won (AUC 0.719 vs 0.707 for $s_1$; 0.691 on dev) with $\tau=0.0122$. The gate is weak: it sends 30\% (practice) and 26\% (dev) of queries to $k=1$, and 27\% and 34\% of those are wrong; a wrong $k=1$ loses the relevant record entirely. At an equal average depth ($k\approx15$) a fixed depth retrieves more (Recall 0.534 vs 0.481 on practice, 0.553 vs 0.503 on dev). On the held-out dev answers the policy scores 0.272 against 0.321 for fixed $k=5$ (difference $-0.049$, CI $-0.123$ to $+0.015$) while using $3.3\times$ the prompt tokens. Even an oracle gate that knows the true rank-1 relevance scores only 0.306 on dev and 0.315 on practice (difference to $k=5$: $-0.015$ and $+0.048$, both not significant). The ceiling is low for two reasons: 13 of the 55 judged queries have no relevant record in the top 20, and even with a relevant record in the prompt 26--45\% of answers score 0 (6/23, 11/34 and 19/42 at $k=1,5,20$), so generation and the judge's gold-only crediting also cap the score. The high arm matters too: on dev, $k=20$ scores below $k=5$ for queries with evidence at ranks 2--20 ($-0.179$, CI $-0.298$ to $-0.060$, $n=7$). In an exploratory replay outside the fixed protocol (which chose $k=20$ on practice, where it replays better: $+0.009$ against $-0.051$ for $k=5$), the same gate with $k\in\{1,5\}$ scores $-0.018$ (CI $-0.068$ to $+0.023$) against fixed $k=5$ on all 55 queries, at $0.80\times$ the tokens.

\subsection{End-to-end cost}

For one spoken question at $k=5$, the stages add up to about 5.4\,s with PyTorch ASR (1555\,ms mean ASR over the 350 practice queries; generation 70\% of the total; encoding and retrieval under 1\%) and about 7.3--7.7\,s with our kernels (per-query ASR scaled by our same-node slowdown: $2.25\times$ on the node that produced the transcripts, $2.45\times$ in the same-job comparison). This is an upper bound, because the 1555\,ms also includes audio loading, feature extraction and process communication, which our kernels do not change, and an estimate, because the stages ran on different machines. It puts the two halves of the work on one scale: our kernels add at most about 1.9--2.3\,s per question, while moving from $k=1$ to $k=20$ adds 2.5\,s (practice) to 5.1\,s (dev) of median generation time and $15\times$ the tokens without a measurable quality gain.

% ============================================================
\section{Discussion and Conclusion}
\label{sec:conclusion}

On the ASR side the evidence is precise: correctness is established at three levels, the cost is measured against the right baseline, and the profile shows that the only kernel worth further work is the matrix multiply. Register tiling (several outputs per thread), vectorised loads and a shape-dependent tile, small for single-row multiplies such as \texttt{lm\_head}, are the next steps. On the RAG side, we detect no quality difference between $k=1$, 5 and 20 in this system under this judge (gains above about 0.11 are excluded, smaller ones are not), so $k=1$ or 5 is the cost-efficient default: it uses a fraction of the tokens. Our own change, adaptive depth, was tested with a protocol fixed in advance and a held-out split, and failed for clear reasons: the confidence signal mis-routes 27--34\% of its confident queries, the $k=20$ arm scores below $k=5$ on dev, and many answers score 0 at every depth.

\textbf{Limitations.} Answer quality rests on 55 judged queries, the first of each split in manifest order rather than a random sample, although their recall matches the full splits (Recall@1 0.249 vs 0.229 and 0.209; Recall@20 0.570 vs 0.579 and 0.580); batches of about 15 queries disagree with each other by as much as any depth effect. The judge is the same model as the generator \citep{zheng2023judging}, credits only the gold records, and cannot tell whether ASR changed the question, because no reference transcripts exist for the queries. Generator and judge ran once per answer; generation variability is not measured. A second judge pass over 56 stored answers (the 36 that scored 0 with a relevant record in the prompt and 20 non-zero ones; three returned unparseable output) gave the same score for 44 of 53 (83\%), with a mean absolute change of 0.05; 4 of the 33 stored zeros rose to 0.25--0.58. So the judge is noisy even at temperature 0; this noise is part of the per-query scores behind our intervals and of the batch differences. Our reading of the 36 zero-score answers finds 12 judge errors (for one query the judge scored a correct answer 0, saying it lacked a detail that it contains; at $k=5$ it scored 4 of 4, and on the second pass 0.58 at $k=1$), 8 answers zeroed for extra content from non-gold records, 14 real misses and 2 unclear cases; this reading is ours, not a measurement. The RAG study used PyTorch-ASR transcripts; our kernels reproduce the test recording's transcript exactly, but we did not compare transcripts for the 475 queries. Generation latency depends on a shared service. The supplied encoder loads in float16 on the CPU; a float32 run gives the same Recall@$k$ to four decimals and the same top record for every query. All GPU results are from one GPU type, and the kernel settings were searched over two constants only. \textbf{What we would do next:} a reranker or query rewriting to reduce retrieval misses, the failure that no choice of depth can fix; and register tiling in the GEMM, the stage that dominates the ASR cost.

% References are excluded from the main-body page limit.
\bibliography{references}
\bibliographystyle{iclr2027_conference}

\appendix
\section{Hand-in output: \texttt{make test} and \texttt{./run.sh --profile}}
\label{app:output}
One Slurm job on one node, with the submitted settings (TILE 32, ATTN\_WARPS 16, ROPE\_THREADS 256, REDUCE\_THREADS 256). The first line is \texttt{nvidia-smi -L}; then the full output of \texttt{make test} and of \texttt{./run.sh --profile} (compiler lines omitted; the home folder is shown as \texttt{<home>}). The same text is \texttt{outputs/handin\_make\_test\_and\_profile.txt} in the supplementary code.
{\fontsize{6.6}{7.7}\selectfont
\begin{verbatim}
GPU 0: NVIDIA GeForce RTX 2080 Ti
GPU: NVIDIA GeForce RTX 2080 Ti (sm_75), 68 SMs, peak memory bandwidth about 616 GB/s

tiny (2 x 3 tokens, Hq=3, Hk=1)      PASS  max_err=1.08e-07     0.004 ms      0.4 GB/s ( 0% of peak)
odd sizes, partial rotation          PASS  max_err=1.08e-07     0.004 ms      3.9 GB/s ( 1% of peak)
audio encoder (partial, 32/64)       PASS  max_err=1.93e-07     0.031 ms    491.8 GB/s (80% of peak)
text decoder (GQA 16/4, full)        PASS  max_err=1.64e-07     0.004 ms    302.4 GB/s (49% of peak)

Task 1 (RoPE): ALL PASS
GPU: NVIDIA GeForce RTX 2080 Ti (sm_75), 68 SMs, peak memory bandwidth about 616 GB/s

rmsnorm, one short row               PASS  max_err=8.27e-08     0.004 ms      0.1 GB/s ( 0% of peak)
rmsnorm, odd cols                    PASS  max_err=3.68e-07     0.004 ms      9.6 GB/s ( 2% of peak)
rmsnorm, text decoder                PASS  max_err=4.35e-07     0.006 ms    164.6 GB/s (27% of peak)
layernorm, odd cols (example)        PASS  max_err=3.91e-07     0.005 ms      8.1 GB/s ( 1% of peak)
layernorm, audio encoder             PASS  max_err=5.42e-07     0.040 ms    381.0 GB/s (62% of peak)
softmax, one short row               PASS  max_err=9.28e-09     0.004 ms      0.0 GB/s ( 0% of peak)
softmax, odd cols                    PASS  max_err=9.75e-10     0.004 ms      5.6 GB/s ( 1% of peak)
softmax, large inputs (+-100)        PASS  max_err=2.67e-08     0.004 ms      4.2 GB/s ( 1% of peak)
softmax, all negative (-200..-150)   PASS  max_err=4.80e-09     0.004 ms      1.2 GB/s ( 0% of peak)
softmax, vocabulary row (59264)      PASS  max_err=6.97e-11     0.118 ms      4.0 GB/s ( 1% of peak)
softmax, 64 x 4096                   PASS  max_err=9.11e-10     0.009 ms    243.0 GB/s (39% of peak)

Task 2 (reductions): ALL PASS
GPU: NVIDIA GeForce RTX 2080 Ti (sm_75), 68 SMs, peak memory bandwidth about 616 GB/s

tiny 3 x 5, K = 7                    PASS  max_err=9.76e-09     0.004 ms        0 GFLOP/s
    for comparison: naive      0 GFLOP/s, cuBLAS      0 GFLOP/s (CPU reference)
odd sizes 37 x 100, K = 33, bias     PASS  max_err=1.55e-07     0.007 ms       35 GFLOP/s
    for comparison: naive     50 GFLOP/s, cuBLAS     21 GFLOP/s (CPU reference)
odd sizes, bias + GELU               PASS  max_err=7.70e-07     0.014 ms       62 GFLOP/s
    for comparison: naive     87 GFLOP/s, cuBLAS     85 GFLOP/s (CPU reference)
decoder q_proj (64x2048x2048)        PASS  max_err=1.10e-05     0.367 ms     1462 GFLOP/s
    for comparison: naive    153 GFLOP/s, cuBLAS   4647 GFLOP/s (cuBLAS reference)
decoder gate_proj (64x6144x2048)     PASS  max_err=1.19e-05     1.076 ms     1497 GFLOP/s
    for comparison: naive    162 GFLOP/s, cuBLAS   6825 GFLOP/s (cuBLAS reference)
decoder down_proj (64x2048x6144)     PASS  max_err=3.62e-05     0.883 ms     1825 GFLOP/s
    for comparison: naive    193 GFLOP/s, cuBLAS   8143 GFLOP/s (cuBLAS reference)
lm_head, last token (1x59264x2048)   PASS  max_err=1.24e-05     3.845 ms       63 GFLOP/s
    for comparison: naive    111 GFLOP/s, cuBLAS    291 GFLOP/s (cuBLAS reference)
encoder q_proj (1500x1280x1280)      PASS  max_err=2.38e-07     2.275 ms     2160 GFLOP/s
    for comparison: naive    231 GFLOP/s, cuBLAS  13201 GFLOP/s (cuBLAS reference)
encoder fc1 + GELU (1500x5120x1280)  PASS  max_err=6.92e-07     9.330 ms     2107 GFLOP/s
    for comparison: naive    233 GFLOP/s, cuBLAS  13363 GFLOP/s (cuBLAS reference)
conv2 via im2col (1500x1280x3840)    PASS  max_err=7.16e-07     6.921 ms     2130 GFLOP/s
    for comparison: naive    232 GFLOP/s, cuBLAS  13340 GFLOP/s (cuBLAS reference)

Task 3 (tiled GEMM): ALL PASS
GPU: NVIDIA GeForce RTX 2080 Ti (sm_75), 68 SMs, peak memory bandwidth about 616 GB/s

tiny, one tile (4 queries, 5 keys)   PASS  max_err=1.11e-07     0.006 ms        0 GFLOP/s
several tiles, peaky scores          PASS  max_err=4.97e-07     0.019 ms       12 GFLOP/s
GQA 4 heads / 2 kv heads             PASS  max_err=1.30e-07     0.008 ms       70 GFLOP/s
causal, seq_q = seq_k                PASS  max_err=4.86e-07     0.022 ms      116 GFLOP/s
causal, 5 new queries, 37 keys       PASS  max_err=1.78e-07     0.010 ms       19 GFLOP/s
causal, one new query                PASS  max_err=3.39e-07     0.019 ms       31 GFLOP/s
text decoder (GQA 16/4, causal)      PASS  max_err=5.21e-07     0.022 ms      779 GFLOP/s
audio encoder (20 heads, 1500)       PASS  max_err=2.68e-07     7.526 ms     1531 GFLOP/s

Task 4 (attention): ALL PASS
GPU: NVIDIA GeForce RTX 2080 Ti (sm_75)
Kernels: rope=cuda  reduce=cuda  gemm=cuda  attention=cuda  provided=cuda
Library: <home>/edin-mls-26-autumn/hw1-asr/build/sm_75/libasr_kernels.so
Loading the model (about 9 GB of float32 weights)...
Loaded in 179.9 s
Audio: 3.50 s, mel features (1, 128, 3000), prompt 59 tokens
  run 1:   2697.3 ms  (13 tokens)
  run 2:   2708.7 ms  (13 tokens)
  run 3:   2728.0 ms  (13 tokens)

GPU time per operation, one transcription (2846.5 ms in kernels):
  operation    runs on    calls   total ms   share
  linear       cuda        2757     2449.5   86.1%
  attention    cuda         396      378.3   13.3%
  rmsnorm      cuda         741        6.7    0.2%
  silu_mul     cuda         364        3.5    0.1%
  rope         cuda         396        3.0    0.1%
  layernorm    cuda          65        2.4    0.1%
  softmax      cuda          13        1.5    0.1%
  embedding    cuda          13        1.3    0.0%
  im2col       cuda           2        0.2    0.0%

Time: 2711.3 ms per transcription (+/- 12.7 ms), 208.6 ms per token
Peak GPU memory: 8.56 GiB
Transcription: Concord returned to its place amidst the tents.
Expected:      Concord returned to its place amidst the tents
Accuracy: 100.0%
Status: PASS
\end{verbatim}
}

\section{Kernel ablations and profiles}
\label{app:ablations}
\begin{table}[h]
\caption{Kernel ablations. Each variant is the same task file built with one or two compile-time flags and checked by the same tests; median of 3 runs of the course test, ms per call, handout settings unless stated.}
\begin{center}
\small
\begin{tabular}{p{3.6cm}p{3.4cm}p{5.6cm}}
\toprule
Kernel / case & Change & Result\\
\midrule
GEMM, encoder fc1 ($1500\times5120\times1280$) & tile 16 $\to$ 32; 8; no $W$ padding & 14.2 $\to$ 9.3\,ms; 24.8; 24.0 (1.7$\times$ slower)\\
GEMM, lm\_head (1 row) & tile 16 $\to$ 8 / 32 & 2.72 $\to$ 1.54 / 4.84\,ms: the best tile depends on the shape\\
Attention, audio encoder & 4 warps $\to$ 1 / 2 / 8 / 16; no $K$ padding & 11.5 $\to$ 103.6 / 30.1 / 8.3 / 7.6\,ms; 54.2\,ms (4.7$\times$)\\
Softmax, vocabulary row & 256 $\to$ 32 / 1024 threads; 3-pass & 0.118 $\to$ 0.886 / 0.041\,ms; 0.150\,ms (+27\%)\\
LayerNorm, 1280-wide row & 256 $\to$ 1024 threads & 1.9$\times$ slower (too little work per thread)\\
RMSNorm, decoder & shuffle $\to$ atomic reduction & 0.006 $\to$ 0.022\,ms (3.7$\times$)\\
GEMM, model shapes & coalesced $\to$ strided $W$ loads & 7--43\% slower (lm\_head 5\%)\\
RoPE, audio encoder & 32--1024 threads; fast math & flat between 64 and 512 (memory-bound, 80\% of peak); fast math fails the test\\
\bottomrule
\end{tabular}
\end{center}
\end{table}

\noindent End to end with all our kernels (one job, one node): GEMM tile 8 4992.2\,ms, tile 16 3288.3\,ms, tile 32 2865.2\,ms. Profile shares on the PyTorch path (same recording): linear 47.6\%, attention 22.5\%, RoPE 25.0\%, RMSNorm 4.2\%; with our kernels: linear 86.1\%, attention 13.3\%, everything else 0.6\%.

\section{Additional RAG detail}
\label{app:rag}
\begin{table}[h]
\caption{Retrieval depth sweep on the practice split ($n=350$, exact search, warm runs): Recall@$k$ with 95\% bootstrap interval, hit@$k$, context size and retrieval time.}
\begin{center}
\small
\begin{tabular}{rrrrr}
\toprule
$k$ & Recall@$k$ (95\% CI) & hit@$k$ & context words & retrieval ms\\
\midrule
1 & 0.229 (0.196--0.265) & 0.411 & 175 & 0.44\\
5 & 0.421 (0.378--0.464) & 0.591 & 897 & 0.45\\
20 & 0.579 (0.534--0.621) & 0.737 & 3640 & 0.44\\
100 & 0.770 (0.732--0.806) & 0.880 & 17957 & 0.45\\
\bottomrule
\end{tabular}
\end{center}
\end{table}
\begin{table}[h]
\caption{Judge sub-scores (0--4 each) on the judged subset, the same queries at every $k$.}
\begin{center}
\small
\begin{tabular}{lccc}
\toprule
correctness / grounding / completeness & $k=1$ & $k=5$ & $k=20$\\
\midrule
practice ($n=28$) & 1.07 / 1.11 / 0.82 & 1.14 / 1.18 / 0.89 & 1.18 / 1.32 / 1.00\\
dev ($n=27$) & 1.15 / 1.30 / 0.96 & 1.37 / 1.44 / 1.04 & 1.04 / 1.04 / 0.96\\
\bottomrule
\end{tabular}
\end{center}
\end{table}

Retrieval: practice split, 350 queries; MRR@10 0.489, NDCG@10 0.420 (dev 0.460 and 0.398). For 12\% of queries no relevant record appears even in the top 100. Of the 55 judged queries, 30, 26 and 31 answers score zero at $k=1$, 5 and 20: the judge credits only the labelled gold records. Judge failures (unparseable output): 2 of 30 practice and 1 of 28 dev queries, excluded at every depth.

\section{Supplementary code}
\label{app:code}
\texttt{hw1-asr/csrc/}: the five task files. \texttt{hw1-asr/ablation/}: variant builder and Slurm jobs. \texttt{rag/extensions/}: transcript cache, depth sweep, answer generation with token counts, paired summary with bootstrap intervals, the post hoc rank-group and replay analyses, the judge test--retest summary, adaptive depth, and the two-stage demo used in the video. \texttt{outputs/}: the hand-in output, the same-job PyTorch comparison and the per-query judge scores, from which every judged-answer and adaptive-policy number can be recomputed. \texttt{EXPERIMENT\_LOG.md}: every run with its command and numbers.

% Keep this line unchanged.
\newpage

\section*{AI declaration}
Parts of this coursework were completed with the assistance of a personal AI assistant (Claude, Anthropic), following the course's AI policy. Claude wrote the CUDA kernels and the experiment scripts under our direction and helped us debug them; it also helped analyse results, keep the experiment log and draft this report, the code README and the narration notes for the project video. We reviewed, tested and edited all of it and can explain every part. All numbers come from runs we executed on the course GPU cluster (speech recognition), a laptop CPU (encoding and retrieval) and the ELM service (generation and judging). We checked every design choice and result and remain responsible for the content. The ELM service (EuroLLM-22B-Instruct) was used only as the system's generator and judge.

\section*{Group information}
% The filled-in page is in group_info.tex, which is git-ignored so that names and UUNs
% never reach the repository. Without that file the template placeholders are printed.
\IfFileExists{group_info.tex}{\input{group_info.tex}}{%
Group ID: XXX

Information of group member A:
\begin{itemize}
    \item Name: John Doe
    \item UUN: sXXXXXXX
\end{itemize}
}

\end{document}
